\section{Conclusion}
\label{sec:conclusion}
In this work, we introduced a Spatio-Spectral Joint Embedding Predictive Architecture (SS-JEPA) specifically designed for the challenges of multimodal Martian landslide segmentation. By shifting from traditional end to end supervised learning methods to a predictive latent-representation framework, our approach effectively overcomes the issues of sensor noise inherent in orbital Martian imagery.

Our primary contribution is a spatio-spectral masking strategy that enables the model to learn deep physical correlations between heterogeneous modalities,  without requiring extensive human-labeled data. Quantitatively, our \(20\)M parameter model achieves an mIoU of \(0.865\), matching \(113\)M  parameters model with a \(\sim 5.6 \times\) reduction in model size while maintaining the state-of-the-art segmentation performance. Beyond these numerical gains, the proposed architecture captures the intrinsic physical corelations between different martian modalities to understand how modalities like topographic steepness and thermal signatures interact to define the Martian landscape.

The learned representations further highlight the importance of cross-modal interactions in multispectral planetary imagery, demonstrating how geological structures can manifest differently across spectral bands while still preserving meaningful latent relationships. Overall, this work presents a robust, scalable, and computationally efficient approach for automated Martian landslide segmentation, providing a practical pathway toward intelligent planetary exploration systems for the Red Planet.


